\chapter{Endomaps, Powers, and Finite Idempotence}
\label{ch05:main}

Chapter~\ref{ch04:main} established that the kernel's objects and arrows
form a category, and observed (Corollary~\ref{ch04:cor:monoid}) that for
each object $X$ the endo-arrows $\Hom(X, X)$ form a monoid under
composition. That monoid is the subject of this chapter, and it is the only
place in the entire kernel where the expression $f^{n}$ --- ``apply $f$,
$n$ times'' --- is granted a meaning. The grant is deliberate and narrow:
powers are defined for a single arrow whose declared source and declared
target are the same object, and for nothing else. Heterogeneous sequences of
arrows, however composable, receive a different construction (the
semigroupoid path of Chapter~\ref{ch06:main}) with deliberately no power
operation, and the type system keeps the two apart.

Within the licensed territory, the mathematics is classical and rich. A
map from a finite set to itself, iterated, exhibits a completely understood
long-run structure: every orbit runs down a finite tail and then cycles
forever, the powers of the map form an eventually periodic sequence, and the
monogenic subsemigroup generated by the map contains a unique idempotent.
The chapter develops this structure theory in full --- functional graphs,
orbit tails and cycles, the index and period of the power sequence, the
cyclic-group kernel of the monogenic subsemigroup --- because it is the
correct setting in which to state and prove the result the kernel actually
ships: for every endomap $f$ on a carrier of size $N$, the power $f^{N!}$
is idempotent,
\[
  f^{N!} \circ f^{N!} = f^{N!}.
\]
We prove this completely (Theorem~\ref{ch05:thm:factorial}), and then ---
because intellectual honesty about bounds is part of the claim discipline
--- we present the sharper statements that the structure theory makes
available: the exact characterization of which powers are idempotent, the
minimal idempotent power in terms of index and period, and the smaller
uniform exponent $\mathrm{lcm}(1, \dots, N)$ that also works. All of this is
standard mathematics, long predating this project
\cite{pin, howie}, and the chapter presents it as such; the kernel's
contribution is the disciplined implementation, and the kernel implements
and tests exactly the factorial bound.

The chapter closes with the computational layer: exponentiation by
squaring, whose correctness we prove against the monoid power laws, and the
bounded presentation-name construction that keeps the \emph{names} of
astronomically large powers small --- bit length plus a short SHA-256 label
--- so that computing $f^{10^{5000}}$ neither materializes a five-thousand-
digit decimal string nor trips the host language's integer-to-string
conversion limit. A short section explains why none of this naming
machinery can affect correctness: names are extensionally inert by
construction, excluded from arrow equality and hashing, and the regression
suite pins the behavior at $2^{30}$ and $10^{5000}$.

\section{The Endomap Construction}
\label{ch05:sec:construction}

\subsection{Definition and the certification reading}

\begin{definition}[Endomap]
\label{ch05:def:endomap}
An \emph{endomap} is an arrow $f : X \to X$ whose declared source and
declared target are the same object --- equal in name and in ordered
carrier, per the object identity of Chapter~\ref{ch03:chap}. The kernel
represents the property by a distinct wrapper type, $\mathrm{Endomap}(f)$,
whose constructor \emph{verifies} $f$'s source-target equality and refuses
construction otherwise, raising the same typed error class used for
composition mismatches. The wrapped arrow's underlying object is called the
\emph{carrier} of the endomap.
\end{definition}

The wrapper deserves a word, because it is easy to misread as bureaucracy.
$\mathrm{Endomap}(f)$ adds no data to $f$: it carries exactly the arrow it
was given. What it adds is a \emph{certificate}, checked once at
construction, that the arrow inhabits an endomorphism monoid --- and
therefore that the operations defined only on such arrows (the powers of
Section~\ref{ch05:sec:powers}) are meaningful for it. Code that receives an
$\mathrm{Endomap}$ value may take powers without re-checking anything; code
that holds a bare arrow may not, and the host language's type distinction
enforces the may-not. In the terminology of typed API design, the wrapper is
a \emph{parse-don't-validate} boundary: the fact ``this arrow is endo'' is
established at one site and thereafter carried by the type rather than
re-derived or, worse, assumed.

Two refusals of the constructor are worth making explicit, because both
correspond to real confusions that arise in informal operator notation.

\begin{ruleenv}[Same carrier is not enough]
\label{ch05:rule:samecarrier}
An arrow between two distinct objects with equal ordered carriers --- say
$f : X \to X''$ where $X$ and $X''$ both carry $(0, 1)$ under different
names --- is \emph{not} an endomap, and $\mathrm{Endomap}(f)$ is refused.
This is the arrow-level consequence of name-inclusive object identity: the
kernel has no basis for asserting that a trip from \texttt{register\_a} to
\texttt{register\_b} can be repeated, because the second application would
need an arrow \emph{out of} \texttt{register\_b}, and $f$ is not one. If the
caller means the two registers to be identified, the identification must be
an explicit arrow, and the composite with it will be an endomap on whichever
object the caller chooses.
\end{ruleenv}

\begin{ruleenv}[Same cardinality is nothing at all]
\label{ch05:rule:samecard}
A fortiori, an arrow between distinct objects of merely equal size is not an
endomap, and no positional re-indexing is performed to make it one. The
untyped habit this blocks is pervasive in matrix-flavored notation, where
any square matrix can be squared; in the kernel, squareness of the index
range does not confer iterability --- only sameness of the object does.
\end{ruleenv}

\subsection{The endomorphism monoid, concretely}

By Corollary~\ref{ch04:cor:monoid}, the endo-arrows on a fixed object $X$
with $|X| = N$ form a monoid of cardinality $N^{N}$: the full transformation
monoid on $N$ points, in its standard-index presentation. It contains the
symmetric group (the $N!$ arrows whose mapping tuples are permutations of
$(0, \dots, N-1)$), the constant maps (the $N$ arrows whose tuples are
constant), the identity $\id_X$, and everything between. Composition of
endomaps is always defined --- the middle object is $X$ on both sides ---
which is precisely why every ordered triple participated in the
associativity sweep of Section~\ref{ch04:sec:assoc-count}.

For the record, and because the number recurs in the verification story
(Section~\ref{ch05:sec:evidence}), the count of \emph{all} endomaps on
carriers of size at most four, one representative object per size, is
\[
  \sum_{k=0}^{4} k^{k} = 1 + 1 + 4 + 27 + 256 = 289,
\]
with the convention $0^{0} = 1$: the empty object carries exactly one
endomap, the empty identity. These $289$ endomaps are the complete
population of the bounded exhaustive idempotence test.

\section{Monoid Power Laws}
\label{ch05:sec:powers}

\begin{definition}[Powers]
\label{ch05:def:power}
Let $f : X \to X$ be an endomap. Define $f^{0} = \id_X$ and, for $n \ge 0$,
$f^{n+1} = f \circ f^{n}$. Every power is again an endomap on $X$, since
composition of endo-arrows on $X$ stays in $\Hom(X, X)$.
\end{definition}

The definition is the usual recursive one; the familiar exponent laws
follow from associativity and neutrality, and we prove them because the
correctness of the fast exponentiation algorithm
(Section~\ref{ch05:sec:binexp}) reduces to them.

\begin{proposition}[Additivity of exponents]
\label{ch05:prop:add}
For all $m, n \ge 0$: $f^{m+n} = f^{m} \circ f^{n}$.
\end{proposition}

\begin{proof}
Induction on $m$, with $n$ fixed. Base case $m = 0$:
$f^{0+n} = f^{n}$ and $f^{0} \circ f^{n} = \id_X \circ f^{n} = f^{n}$ by
left neutrality (Theorem~\ref{ch04:thm:laws}, Clause~4). Inductive step:
assume $f^{m+n} = f^{m} \circ f^{n}$. Then
\[
  f^{(m+1)+n}
  = f^{(m+n)+1}
  = f \circ f^{m+n}
  = f \circ (f^{m} \circ f^{n})
  = (f \circ f^{m}) \circ f^{n}
  = f^{m+1} \circ f^{n},
\]
using the definition, the inductive hypothesis, and associativity
(Clause~5).
\end{proof}

\begin{corollary}[Powers commute with each other]
\label{ch05:cor:commute}
For all $m, n \ge 0$: $f^{m} \circ f^{n} = f^{n} \circ f^{m}$.
\end{corollary}

\begin{proof}
Both sides equal $f^{m+n}$ by Proposition~\ref{ch05:prop:add} and the
commutativity of addition on natural numbers.
\end{proof}

\begin{proposition}[Multiplicativity of exponents]
\label{ch05:prop:mult}
For all $m, n \ge 0$: $(f^{m})^{n} = f^{mn}$.
\end{proposition}

\begin{proof}
Induction on $n$. Base case: $(f^{m})^{0} = \id_X = f^{0} = f^{m \cdot 0}$.
Inductive step: assuming $(f^{m})^{n} = f^{mn}$,
\[
  (f^{m})^{n+1}
  = f^{m} \circ (f^{m})^{n}
  = f^{m} \circ f^{mn}
  = f^{m + mn}
  = f^{m(n+1)},
\]
using Proposition~\ref{ch05:prop:add} in the third step.
\end{proof}

\begin{remark}[The monogenic submonoid]
\label{ch05:rem:monogenic}
Propositions~\ref{ch05:prop:add} and~\ref{ch05:prop:mult} say that the map
$n \mapsto f^{n}$ is a monoid homomorphism from $(\mathbb{N}, +, 0)$ to
$(\Hom(X, X), \circ, \id_X)$. Its image
$\{ \id_X, f, f^{2}, f^{3}, \dots \}$ is the \emph{monogenic} (single-
generator) submonoid generated by $f$; discarding the identity gives the
monogenic subsemigroup $\langle f \rangle = \{ f, f^{2}, f^{3}, \dots \}$
in the sense standard in semigroup theory \cite{howie}. Because
$\Hom(X, X)$ is finite, this image is finite, so the sequence of powers
must repeat --- and the precise pattern of repetition is the structure
theory of the next section.
\end{remark}

\section{Functional-Graph Structure Theory}
\label{ch05:sec:structure}

Everything provable about iteration of a single finite endomap flows from
one picture: the \emph{functional graph}. This section develops the picture
and its consequences with full proofs. The material is classical --- it is
the elementary structure theory of finite transformation semigroups and of
monogenic semigroups, for which we cite Howie \cite{howie} and Pin
\cite{pin} as standard references --- and we present it as background, with
the kernel-relevant statements extracted as numbered results.

\subsection{The functional graph and its shape}

\begin{definition}[Functional graph]
\label{ch05:def:funcgraph}
Let $f : X \to X$ be an endomap with $|X| = N$, acting on standard indices.
The \emph{functional graph} of $f$ is the directed graph $G_f$ on vertex
set $\{0, \dots, N-1\}$ with an edge $x \to f(x)$ for every vertex $x$.
Every vertex has out-degree exactly one; in-degree is unconstrained (it is
the fiber size $|f^{-1}(x)|$).
\end{definition}

\begin{definition}[Orbit, tail, cycle]
\label{ch05:def:orbit}
The \emph{orbit} of $x$ under $f$ is the sequence
$x, f(x), f^{2}(x), f^{3}(x), \dots$. A vertex $x$ is \emph{cyclic} (or
periodic) if $f^{k}(x) = x$ for some $k \ge 1$; the set of cyclic vertices
is denoted $\mathrm{Cyc}(f)$. The \emph{tail length} of $x$, written
$t(x)$, is the least $j \ge 0$ such that $f^{j}(x)$ is cyclic (we show in
Lemma~\ref{ch05:lem:enters} that such $j$ exists). For a cyclic vertex $c$,
its \emph{cycle length} is the least $k \ge 1$ with $f^{k}(c) = c$.
\end{definition}

\begin{lemma}[Every orbit repeats within $N$ steps]
\label{ch05:lem:repeat}
For every $x$, there exist $0 \le a < b \le N$ with $f^{a}(x) = f^{b}(x)$.
\end{lemma}

\begin{proof}
The $N+1$ values $f^{0}(x), f^{1}(x), \dots, f^{N}(x)$ all lie in the
$N$-element vertex set, so by the pigeonhole principle two of them coincide;
take $a < b$ to be the positions of a coincidence.
\end{proof}

\begin{lemma}[Every orbit enters a cycle; tail and cycle bounds]
\label{ch05:lem:enters}
For every $x$, the tail length $t(x)$ is well defined, and if $\ell(x)$
denotes the cycle length of the cyclic vertex $f^{t(x)}(x)$, then
\[
  t(x) + \ell(x) \le N,
  \qquad\text{hence}\qquad
  t(x) \le N - 1
  \quad\text{and}\quad
  \ell(x) \le N .
\]
\end{lemma}

\begin{proof}
Take $a < b$ as in Lemma~\ref{ch05:lem:repeat} with $b - a$ minimal among
coincidences with minimal $a$; write $y = f^{a}(x)$ and $k = b - a \ge 1$.
Then $f^{k}(y) = f^{b}(x) = f^{a}(x) = y$, so $y$ is cyclic and $t(x) \le a$
is well defined. Moreover the vertices
$f^{0}(x), f^{1}(x), \dots, f^{a + k - 1}(x)$ are pairwise distinct: a
coincidence among them would contradict either the minimality of $a$ (if it
involved an index below $a$) or the minimality of $k$ (if both indices were
at least $a$; two equal values on the cycle at distance less than $k$ would
shorten the cycle). Hence $a + k \le N$. Since $t(x) \le a$ and the cycle
containing $f^{t(x)}(x)$ is the same cycle of length $\ell(x) = k$
(once the orbit meets a cyclic vertex it stays on that vertex's cycle
forever, as $f$ maps each cycle to itself), we get
$t(x) + \ell(x) \le a + k \le N$. The two displayed consequences follow
from $\ell(x) \ge 1$ and $t(x) \ge 0$.
\end{proof}

\begin{proposition}[Shape of the functional graph]
\label{ch05:prop:shape}
Each weakly connected component of $G_f$ contains exactly one cycle, and the
component consists of that cycle together with disjoint trees rooted on the
cycle's vertices, all tree edges oriented toward the cycle. The cyclic
vertices are exactly the vertices of these cycles, and
$\mathrm{Cyc}(f) = \im(f^{N})$: the cyclic vertices are precisely the
eventual image of $f$.
\end{proposition}

\begin{proof}
Existence of at least one cycle per component: pick any vertex of the
component and follow out-edges; by Lemma~\ref{ch05:lem:enters} the walk
enters a cycle, which lies in the same component. Uniqueness: distinct
cycles are vertex-disjoint (a shared vertex would have two distinct
out-edges along the two cycles, contradicting out-degree one), and if a
component contained two disjoint cycles, consider any undirected path
between them; some vertex on that path would need out-degree two or an
out-edge leaving both cycles' basins consistently --- more carefully, every
vertex's forward orbit terminates on exactly one cycle (out-degree one makes
the forward walk deterministic), so ``which cycle does the forward orbit
reach'' is a well-defined function on vertices, constant along every edge in
both directions ($x$ and $f(x)$ reach the same cycle); hence it is constant
on each weakly connected component, so a component sees only one cycle.
Removing the unique cycle's edges leaves each remaining vertex with its
single out-edge pointing one step closer (in tail length) to the cycle,
which is exactly a forest of trees rooted on cycle vertices, oriented
rootward.

For the last claim: if $c$ is cyclic with cycle length $\ell$, then $c =
f^{\ell m}(c)$ for all $m$, and choosing $\ell m \ge N$ exhibits $c \in
\im(f^{N})$. Conversely if $y = f^{N}(x)$ for some $x$, then since
$t(x) \le N - 1 < N$ by Lemma~\ref{ch05:lem:enters}, the orbit of $x$ has
already entered its cycle by step $N$, so $y$ is cyclic.
\end{proof}

\begin{remark}
\label{ch05:rem:permutation}
On the eventual image $\mathrm{Cyc}(f)$, the restriction of $f$ is a
bijection: it permutes each cycle by rotation. All the ``loss'' a finite
endomap can inflict happens in finitely many steps, after which the dynamics
is a permutation of a (possibly much smaller) invariant set. This
two-phase structure --- transient contraction, then permutation --- is the
conceptual content behind every bound in this chapter.
\end{remark}

\subsection{Index and period of the power sequence}

We now shift attention from orbits of points to the sequence of arrows
$f, f^{2}, f^{3}, \dots$ itself, which lives in the finite monoid
$\Hom(X, X)$ and is therefore eventually periodic. The right vocabulary is
the index and period of an eventually periodic sequence, standard in the
theory of monogenic semigroups \cite{howie, pin}.

\begin{definition}[Index and period]
\label{ch05:def:indexperiod}
The \emph{index} $i(f)$ is the least $i \ge 0$ such that $f^{i} = f^{i + k}$
for some $k \ge 1$. The \emph{period} $p(f)$ is the least $k \ge 1$ such
that $f^{i(f)} = f^{i(f) + k}$.
\end{definition}

(With the monoid convention $f^{0} = \id_X$, a permutation has index $0$;
readers of semigroup-theoretic sources, where exponents start at $1$, should
note the harmless offset.)

\begin{proposition}[Structure of the power sequence]
\label{ch05:prop:powerseq}
Write $i = i(f)$ and $p = p(f)$. Then:
\begin{enumerate}
\item The arrows $f^{0}, f^{1}, \dots, f^{i + p - 1}$ are pairwise
distinct, and for all $n \ge i$, $f^{n + p} = f^{n}$.
\item For $m, n \ge i$: $f^{m} = f^{n}$ if and only if
$m \equiv n \pmod{p}$.
\item For $m < i$ and any $n \ne m$: $f^{m} \ne f^{n}$.
\item $i(f) = \max_{x} t(x)$ and $p(f) = \mathrm{lcm}\,
\{\, \ell(c) : c \in \mathrm{Cyc}(f) \,\}$, the least common multiple of
the cycle lengths.
\item $i(f) \le N - 1$ when $N \ge 1$ (and $i(f) = 0$ when $N = 0$), and
$p(f) \le$ the maximal order of a permutation of at most $N$ points; in
particular every cycle length divides $p(f)$ and $p(f)$ divides
$\mathrm{lcm}(1, 2, \dots, N)$ for $N \ge 1$.
\end{enumerate}
\end{proposition}

\begin{proof}
(1) By definition of $i$ there is $k \ge 1$ with $f^{i} = f^{i+k}$, and $p$
is the least such $k$. Composing $f^{i} = f^{i+p}$ with $f^{n-i}$ on the
appropriate side (Proposition~\ref{ch05:prop:add}) gives $f^{n} = f^{n+p}$
for all $n \ge i$. Distinctness: a coincidence $f^{a} = f^{b}$ with
$a < b \le i + p - 1$ would either have $a < i$, contradicting minimality
of $i$ (since then $f^{a}$ recurs), or $a \ge i$ with $b - a < p$,
contradicting minimality of $p$ after translating down to exponent $i$
(compose both sides with $f^{i - a}$... more precisely, from $f^{a} =
f^{b}$ and (1)'s recurrence one derives $f^{i} = f^{i + (b-a)}$ by adding
$f^{p}$-multiples until both exponents exceed $i$ and reducing; the least-
period property then forces $p \mid b - a$, impossible for
$0 < b - a < p$).

(2) The ``if'' direction is (1) iterated. For ``only if,'' suppose
$m, n \ge i$, $m < n$, $f^{m} = f^{n}$, and let $d = n - m$. Adding
multiples of $p$ to both exponents using (1), we may arrange comparisons at
exponent exactly $i$: $f^{i} = f^{i + d'}$ where $d' \equiv d \pmod p$ and
$0 \le d' < p$; minimality of $p$ forces $d' = 0$, i.e.\ $p \mid d$.

(3) If $m < i$ and $f^{m} = f^{n}$ with $n > m$, then $f^{m}$ recurs in the
sequence, and composing the relation forward shows $f^{m} = f^{m + (n-m)s}$
for all $s \ge 1$, exhibiting an index at most $m < i$ --- contradiction.

(4) For the index: unwind what $f^{a} = f^{b}$ means pointwise. For every
$x$, once $a \ge t(x)$, the value $f^{a}(x)$ is on $x$'s cycle and
thereafter the sequence $a \mapsto f^{a}(x)$ is purely periodic with period
$\ell(x)$. If $a \ge \max_x t(x)$ then \emph{every} point is on-cycle, and
taking $k = \mathrm{lcm}$ of all cycle lengths gives $f^{a+k} = f^{a}$;
hence $i(f) \le \max_x t(x)$. Conversely, if $a < t(x_0)$ for some $x_0$,
then $f^{a}(x_0)$ is a tail vertex; tail vertices are never revisited
(their in-orbit occurrence is unique since the orbit is injective until it
meets the cycle, by the distinctness argument in
Lemma~\ref{ch05:lem:enters}), so $f^{b}(x_0) \ne f^{a}(x_0)$ for all
$b > a$, and $f^{a}$ never recurs; hence $i(f) \ge t(x_0)$ for the
maximizing $x_0$. For the period: with all points on-cycle at exponent
$i$, $f^{i+k} = f^{i}$ holds iff $k$ is a multiple of every cycle length
--- rotation of a cycle of length $\ell$ by $k$ steps is the identity on
that cycle iff $\ell \mid k$ --- so the least valid $k$ is the lcm.

(5) The index bound is Lemma~\ref{ch05:lem:enters} ($t(x) \le N - 1$)
combined with (4). For the period: by (4), $p(f)$ is the lcm of the cycle
lengths, which are the cycle lengths of the permutation $f$ restricted to
$\mathrm{Cyc}(f)$, a set of at most $N$ points; the lcm of the lengths of
disjoint cycles fitting in $N$ points is at most the maximal order of a
permutation on $N$ points (Landau's function), and since each individual
cycle length is at most $N$, the lcm divides $\mathrm{lcm}(1, \dots, N)$.
\end{proof}

\subsection{The monogenic subsemigroup and its unique idempotent}

\begin{proposition}[Structure of $\langle f \rangle$; standard \cite{howie, pin}]
\label{ch05:prop:monogenic}
Let $i = i(f)$, $p = p(f)$, and let
$\langle f \rangle = \{ f, f^{2}, f^{3}, \dots \}$ be the monogenic
subsemigroup. Write $i' = \max(i, 1)$. Then:
\begin{enumerate}
\item $\langle f \rangle = \{ f, f^{2}, \dots, f^{i' + p - 1} \}$, with
these elements pairwise distinct; in particular
$|\langle f \rangle| = i' + p - 1$.
\item The subset $K_f = \{ f^{i'}, f^{i' + 1}, \dots, f^{i' + p - 1} \}$ is
closed under composition and is a cyclic group of order $p$.
\item The identity element of $K_f$ is $f^{k^{*}}$, where $k^{*}$ is the
unique integer with $i' \le k^{*} \le i' + p - 1$ and
$p \mid k^{*}$; equivalently, $k^{*}$ is the least positive multiple of $p$
that is at least $i$.
\item $f^{k^{*}}$ is the unique idempotent element of $\langle f \rangle$.
\end{enumerate}
\end{proposition}

\begin{proof}
(1) is Proposition~\ref{ch05:prop:powerseq}(1)--(3) restricted to positive
exponents. (2) Closure: for $a, b \ge i'$, $f^{a} \circ f^{b} = f^{a+b}$
with $a + b \ge i'$, and by
Proposition~\ref{ch05:prop:powerseq}(2) this equals the unique
representative $f^{c}$ with $i' \le c \le i' + p - 1$ and
$c \equiv a + b \pmod p$. Thus $K_f$ with composition is isomorphic to the
additive group $\mathbb{Z}/p\mathbb{Z}$ via $f^{a} \mapsto a \bmod p$
(a bijection on $K_f$ by (1)-distinctness and (2)); a semigroup isomorphic
to a group is a group, cyclic of order $p$, generated by $f^{i'}$'s
appropriate power. (3) The identity of $K_f$ corresponds to residue $0$,
i.e.\ the unique $k^{*}$ in the stated window divisible by $p$; the window
$[\,i', i' + p - 1\,]$ contains exactly one multiple of $p$, and it is the
least positive multiple of $p$ that is $\ge i$ (when $i = 0$ this is $p$
itself, matching $i' = 1$). (4) An idempotent $f^{k}$ satisfies
$f^{2k} = f^{k}$. If $k < i$ this contradicts
Proposition~\ref{ch05:prop:powerseq}(3); so $k \ge \max(i,1) = i'$, and by
(2), $2k \equiv k \pmod p$, i.e.\ $p \mid k$. Hence $f^{k} \in K_f$ with
residue $0$, so $f^{k} = f^{k^{*}}$. Conversely $f^{k^{*}}$ is idempotent
since $k^{*} \ge i$ and $p \mid k^{*}$ give
$f^{2k^{*}} = f^{k^{*}}$ by Proposition~\ref{ch05:prop:powerseq}(2).
\end{proof}

\begin{remark}[Terminology across the literature]
\label{ch05:rem:terminology}
The vocabulary of this section varies across sources, and since the
monograph cites the semigroup literature as background, a small
concordance prevents misreadings. What we call the index and period of
$f$ are the standard parameters of a monogenic (equivalently: cyclic)
semigroup, often written $m$ and $r$, with the pair called the
\emph{type} of the semigroup \cite{howie}; the cyclic group $K_f$ of
Proposition~\ref{ch05:prop:monogenic} is the \emph{kernel} of the
monogenic semigroup in the semigroup-theoretic sense --- an unfortunate
but entrenched collision with the linear-algebraic ``kernel'' and with
this monograph's use of ``kernel'' for the software artifact, which is
why we consistently write $K_f$ and never overload the word. The
functional-graph picture, with its rho-shaped orbits, is the standard
combinatorial companion of these parameters and underlies both the
cycle-detection algorithms of computational practice and the finite-
semigroup structure theory surveyed by Pin \cite{pin}. Nothing in this
section is original; the presentation is arranged so that each numbered
statement is exactly what a later chapter or fixture cites.
\end{remark}

This completes the background. Every finite endomap generates a finite
monogenic semigroup whose tail has length governed by the longest orbit
tail and whose head is a cyclic group whose order is the lcm of the cycle
lengths; the unique idempotent sits at the first exponent that is
simultaneously past the tail and divisible by the period. The theorem the
kernel ships is now a corollary of choosing one crude, uniform exponent
that provably satisfies both conditions.

\subsection{Worked examples on four points}
\label{ch05:sec:worked}

Three concrete endomaps on a four-element carrier
$X = (\texttt{state}, (0, 1, 2, 3))$ exercise every feature of the
structure theory, and we work each one completely; the reader is
encouraged to draw the three functional graphs, which make every claim
visible at a glance.

\begin{example}[Mixed tail and cycle]
\label{ch05:ex:mixed}
Let $f$ have mapping tuple $(1, 2, 1, 0)$: so $0 \mapsto 1$,
$1 \mapsto 2$, $2 \mapsto 1$, $3 \mapsto 0$. The orbit of $1$ is
$1 \to 2 \to 1$: a $2$-cycle, so $\mathrm{Cyc}(f) = \{1, 2\}$ and both
cyclic points have $t = 0$, $\ell = 2$. The orbit of $0$ is
$0 \to 1 \to 2 \to 1 \cdots$: tail length $t(0) = 1$. The orbit of $3$ is
$3 \to 0 \to 1 \to \cdots$: tail length $t(3) = 2$. Hence by
Proposition~\ref{ch05:prop:powerseq}(4),
\[
  i(f) = \max(0, 0, 1, 2) = 2,
  \qquad
  p(f) = \mathrm{lcm}\{2\} = 2,
\]
and by Proposition~\ref{ch05:prop:whichpowers} the minimal idempotent
power is $k^{*} = 2$ (the least positive multiple of $2$ that is at least
$2$). Check directly: $f^{2}$ has tuple
$\bigl(f(1), f(2), f(1), f(0)\bigr) = (2, 1, 2, 1)$, and composing
$f^{2}$ with itself indexes $(2,1,2,1)$ at its own entries:
$\bigl(f^{2}(2), f^{2}(1), f^{2}(2), f^{2}(1)\bigr) = (2, 1, 2, 1)$ ---
idempotent, as predicted. The safe exponent of
Theorem~\ref{ch05:thm:factorial} is $4! = 24$, and indeed
$f^{24} = f^{2}$ because $24 \ge i(f) = 2$ and $24 \equiv 2 \pmod 2$
(Proposition~\ref{ch05:prop:powerseq}(2)), confirming
Corollary~\ref{ch05:cor:whichidem}: the factorial power lands exactly on
the monogenic semigroup's unique idempotent, here reached already at
exponent $2$. The gap between $24$ and $2$ is the price of uniformity
discussed in Remark~\ref{ch05:rem:whyfactorial}.
\end{example}

\begin{example}[Pure permutation]
\label{ch05:ex:perm}
Let $\sigma$ have tuple $(1, 2, 3, 0)$: a single $4$-cycle. Every point is
cyclic, so $i(\sigma) = 0$ and $p(\sigma) = 4$; the minimal idempotent
power is $k^{*} = p = 4$, and $\sigma^{4} = \id_X$ --- for a permutation,
the unique idempotent in the monogenic semigroup is the identity, reached
at the permutation's order. The factorial bound gives
$\sigma^{24} = (\sigma^{4})^{6} = \id_X$ by
Proposition~\ref{ch05:prop:mult}, idempotent trivially. Note how the two
proof obligations of Remark~\ref{ch05:rem:whyfactorial} split here: the
index condition is vacuous ($i = 0$) and the entire content is
divisibility of the exponent by the cycle length.
\end{example}

\begin{example}[Constant map]
\label{ch05:ex:const}
Let $c$ have tuple $(2, 2, 2, 2)$. The point $2$ is a fixed point
($1$-cycle); every other point has tail length $1$. So $i(c) = 1$,
$p(c) = 1$, $k^{*} = 1$: the map is already idempotent, $c \circ c = c$,
as direct computation confirms ($c(c(x)) = c(2) = 2 = c(x)$). Here the
divisibility condition is vacuous and the entire content is the index. The
monogenic subsemigroup is the singleton $\{c\}$, the degenerate case
$i' = 1$, $p = 1$ of Proposition~\ref{ch05:prop:monogenic}.
\end{example}

\subsection{Idempotents as retractions}
\label{ch05:sec:retractions}

The idempotent that the factorial power produces has a clean pointwise
characterization worth recording, both because it explains \emph{what} the
stabilized map is and because it connects the algebra back to the
functional-graph picture.

\begin{proposition}[Idempotent endomaps are retractions onto their fixed
points]
\label{ch05:prop:retraction}
Let $e : X \to X$ be an endomap. The following are equivalent:
\begin{enumerate}
\item $e$ is idempotent, $e \circ e = e$;
\item every point of $\im(e)$ is a fixed point of $e$;
\item $e$ acts as the identity on its image.
\end{enumerate}
In that case $\im(e) = \mathrm{Fix}(e)$, and $e$ is a retraction of $X$
onto its fixed-point set: it maps every point somewhere, and points
already in the image stay put.
\end{proposition}

\begin{proof}
(1) $\Rightarrow$ (2): let $y \in \im(e)$, say $y = e(x)$; then
$e(y) = e(e(x)) = e(x) = y$. (2) $\Leftrightarrow$ (3) is a rephrasing.
(2) $\Rightarrow$ (1): for any $x$, $e(x) \in \im(e)$, so
$e(e(x)) = e(x)$; pointwise agreement at every index plus shared declared
types gives $e \circ e = e$. For the final claim:
(2) gives $\im(e) \subseteq \mathrm{Fix}(e)$, and any fixed point $y$
satisfies $y = e(y) \in \im(e)$, so the two sets coincide.
\end{proof}

For the factorial power specifically, the fixed-point set of $f^{N!}$ is
exactly the cyclic vertex set $\mathrm{Cyc}(f)$: every cyclic vertex is
returned to itself by an exponent divisible by its cycle length, and every
tail vertex has left the tail by exponent $N!$ and cannot be its own
image. So Theorem~\ref{ch05:thm:factorial} has a picture: \emph{the
factorial power is the retraction of the functional graph onto its
cycles}, collapsing each tree to the cycle vertex its flow reaches, and
freezing the cycles pointwise. In Example~\ref{ch05:ex:mixed}, $f^{2}$
sends $0, 2 \mapsto 2$ and $1, 3 \mapsto 1$: the four-point graph retracts
onto the $2$-cycle $\{1, 2\}$, exactly as the picture demands.

\section{The Factorial Idempotence Theorem}
\label{ch05:sec:factorial}

\begin{theorem}[Safe finite power]
\label{ch05:thm:factorial}
Let $f : X \to X$ be an endomap on a carrier of size $N = |X|$. Then
$f^{N!}$ is idempotent:
\[
  f^{N!} \circ f^{N!} = f^{N!}.
\]
\textup{Status: standard result \cite{pin, howie}; complete informal proof
below; verified exhaustively by the executable suite for all $289$ endomaps
on carriers of size at most four; Lean formalization open
(ledger PO-TKS8-CORE-04).}
\end{theorem}

We give two proofs: a self-contained pointwise proof, which is the argument
the seed paper records and which requires nothing beyond
Lemma~\ref{ch05:lem:enters}, and a one-line derivation from the structure
theory, which situates the theorem correctly as the crude-bound corollary
of Proposition~\ref{ch05:prop:monogenic}.

\begin{proof}[First proof (self-contained, pointwise)]
The cases $N = 0$ and $N = 1$ are immediate: for $N = 0$ the only endomap
is the empty identity and every power equals it; for $N = 1$ the only
endomap is $\id_X$ and likewise. Assume $N \ge 2$, and fix an arbitrary
point $x$. By Lemma~\ref{ch05:lem:enters}, the tail length satisfies
$t(x) \le N - 1$. Since $N \ge 2$ gives $N! \ge N > N - 1 \ge t(x)$, the
point $f^{N!}(x)$ lies on $x$'s cycle. Let $\ell \le N$ be that cycle's
length (Lemma~\ref{ch05:lem:enters} again). Because $\ell \le N$, the
integer $\ell$ divides $N!$ --- it is one of the factors $1, 2, \dots, N$
of the product --- so a further $N!$ applications of $f$ traverse the cycle
a whole number $N!/\ell$ of times and return to the same vertex:
\[
  f^{2 N!}(x) = f^{N!}\bigl(f^{N!}(x)\bigr) = f^{N!}(x).
\]
The point $x$ was arbitrary, so the arrows $f^{2N!}$ and $f^{N!}$ agree at
every standard index; they share declared source and target $X$; hence
$f^{2N!} = f^{N!}$ as arrows, and by
Proposition~\ref{ch05:prop:add} this is exactly
$f^{N!} \circ f^{N!} = f^{N!}$.
\end{proof}

\begin{proof}[Second proof (via structure theory)]
By Proposition~\ref{ch05:prop:powerseq}(4,5), $i(f) \le \max(N-1, 0) \le
N!$ and $p(f)$ is the lcm of cycle lengths, each at most $N$, so
$p(f) \mid N!$. Thus $N! \ge i(f)$ and $p(f) \mid N!$, and
Proposition~\ref{ch05:prop:powerseq}(2) applied to the exponents $N!$ and
$2N!$ (both $\ge i(f)$, congruent mod $p(f)$) gives $f^{2N!} = f^{N!}$.
(For $N \in \{0, 1\}$, $N! = 1$ and every endomap on at most one point is
the identity, idempotent trivially --- note the $N=1$ case needs the
separate sentence because $1! = 1$ and $f^{1} = f$, which is idempotent
only because $f = \id_X$ there.)
\end{proof}

\begin{remark}[Why the factorial]
\label{ch05:rem:whyfactorial}
The exponent $N!$ is doing two jobs at once, and the proof shows them
cleanly separated: it must be \emph{large enough} to outrun every tail
($N! \ge N - 1$, with room to spare), and it must be \emph{divisible
enough} to be a multiple of every possible cycle length
($\ell \mid N!$ for all $\ell \le N$). Any exponent with both properties
works; $N!$ is the crudest natural choice with a one-line proof of each
property. Its virtue is not sharpness but \emph{uniformity and
checkability}: it depends only on the carrier size, not on any structural
analysis of $f$, so the implementation can compute the safe exponent
without first computing the functional graph. The next section quantifies
what sharper choices exist and what they would cost.
\end{remark}

\begin{corollary}[The factorial power is the idempotent of $\langle f \rangle$
--- when $N \ge 2$]
\label{ch05:cor:whichidem}
For $N \ge 2$, $f^{N!}$ equals $f^{k^{*}}$, the unique idempotent of the
monogenic subsemigroup of Proposition~\ref{ch05:prop:monogenic}.
\end{corollary}

\begin{proof}
$f^{N!}$ is an idempotent element of $\langle f \rangle$ (it is a positive
power of $f$, since $N! \ge 1$, and idempotent by
Theorem~\ref{ch05:thm:factorial}); uniqueness in
Proposition~\ref{ch05:prop:monogenic}(4) forces
$f^{N!} = f^{k^{*}}$.
\end{proof}

\section{Sharper Variants, Stated Honestly}
\label{ch05:sec:sharper}

The factorial bound invites two immediate refinements, both classical, and
this section presents them with the same care as the theorem itself ---
including the explicit statement of what the kernel does and does not
implement.

\subsection{Exact characterization of idempotent powers}

The structure theory already delivered the exact answer, and we restate it
in bound-free form:

\begin{proposition}[Which powers are idempotent]
\label{ch05:prop:whichpowers}
For $k \ge 1$, the power $f^{k}$ is idempotent if and only if
$k \ge i(f)$ and $p(f) \mid k$. The minimal such $k$ is
\[
  k^{*}(f) =
  \begin{cases}
    p(f) \cdot \bigl\lceil i(f) / p(f) \bigr\rceil, & i(f) \ge 1, \\[2pt]
    p(f), & i(f) = 0,
  \end{cases}
\]
that is, the least positive multiple of the period that is at least the
index.
\end{proposition}

\begin{proof}
Contained in the proof of Proposition~\ref{ch05:prop:monogenic}(3,4).
\end{proof}

The minimal exponent $k^{*}(f)$ can be dramatically smaller than $N!$: for
a permutation consisting of a single $N$-cycle, $i = 0$ and $p = N$, so
$k^{*} = N$ versus $N!$; for a map with one long tail into a fixed point,
$p = 1$ and $k^{*} = i \le N - 1$. But computing $k^{*}(f)$ requires
knowing $i(f)$ and $p(f)$, i.e.\ analyzing the functional graph of the
specific $f$ --- a linear-time computation, but a computation about $f$
rather than about $N$ alone.

\subsection{The lcm bound for the uniform exponent}

Between the exact per-map answer and the factorial there is a natural
intermediate: keep the exponent uniform in $N$, but replace $N!$ by the
least common multiple of $1$ through $N$.

\begin{proposition}[Uniform lcm exponent]
\label{ch05:prop:lcm}
Let $L(N) = \mathrm{lcm}(1, 2, \dots, N)$ for $N \ge 1$. For every endomap
$f$ on a carrier of size $N \ge 1$, the power $f^{L(N)}$ is idempotent.
\end{proposition}

\begin{proof}
By Proposition~\ref{ch05:prop:powerseq}(5), $p(f) \mid L(N)$. For the index:
$i(f) \le N - 1$, and $L(N) \ge N > N - 1$ for $N \ge 1$ (indeed
$N \mid L(N)$ so $L(N) \ge N$). Hence $L(N) \ge i(f)$ and
$p(f) \mid L(N)$, and Proposition~\ref{ch05:prop:whichpowers} applies.
(For $N = 1$, $L(1) = 1$ and the unique endomap is the identity.)
\end{proof}

The gap between the two uniform exponents is enormous: $L(N)$ grows like
$e^{(1 + o(1)) N}$ (by the prime number theorem, since $\log L(N)$ is the
Chebyshev function $\psi(N)$), while $N!$ grows like $e^{N \log N - N}$
by Stirling. Already at $N = 10$: $10! = 3{,}628{,}800$ versus
$L(10) = 2{,}520$. Both are standard facts of elementary number theory
combined with the semigroup structure theory above \cite{pin, howie}; the
maximal possible \emph{period} over all endomaps on $N$ points is Landau's
function (the maximal order of a permutation of at most $N$ points), which
is smaller still, but the exponent must also clear the index, so the clean
uniform statements are the two just given.

The divergence is visible immediately in a small table:

\begin{center}
\begin{tabular}{@{}rrr@{}}
\toprule
$N$ & $N!$ & $L(N) = \mathrm{lcm}(1, \dots, N)$ \\
\midrule
$1$ & $1$ & $1$ \\
$2$ & $2$ & $2$ \\
$3$ & $6$ & $6$ \\
$4$ & $24$ & $12$ \\
$5$ & $120$ & $60$ \\
$6$ & $720$ & $60$ \\
$7$ & $5{,}040$ & $420$ \\
$8$ & $40{,}320$ & $840$ \\
\bottomrule
\end{tabular}
\end{center}

The two agree through $N = 3$ and separate permanently at $N = 4$, where
the factorial first multiplies in a factor ($4$) whose prime content
($2^2$) is only partially new. Neither column is remotely tight for a
\emph{typical} endomap --- Example~\ref{ch05:ex:mixed}'s map on four
points stabilizes at exponent $2$ against the factorial's $24$ --- but
tightness is not the design goal; a uniform, structure-blind, provably
safe exponent is, and for that role either column serves, with $N!$
chosen for the shorter proof.

\subsection{What the kernel implements}

\begin{finding}
\label{ch05:find:implemented}
The kernel implements and tests \emph{the factorial bound and only the
factorial bound}. The endomap type exposes an operation returning
$f^{|X|!}$ (computed by binary exponentiation,
Section~\ref{ch05:sec:binexp}), together with a predicate testing
idempotence of an arbitrary endomap by direct comparison of $f \circ f$
with $f$. It does not compute $i(f)$, $p(f)$, $k^{*}(f)$, or $L(N)$, and
no kernel API claims minimality of the factorial exponent. The sharper
variants of this section are presented as standard mathematics
\cite{pin, howie} for the reader's orientation and as candidate future
obligations; they carry no executable or formal artifact in this kernel,
and nothing downstream may cite them as implemented.
\end{finding}

The design rationale for shipping only the crude bound is the same
throughout the kernel: prefer the statement whose proof is short, whose
implementation is a one-liner over an already-tested primitive, and whose
correctness does not depend on a per-instance structural analysis that
would itself need tests. A minimal-exponent API would be genuinely useful
--- and genuinely riskier, since an off-by-one in tail computation would
produce non-idempotent ``idempotents'' silently. The kernel takes the safe
exponent; the analysis layer can be built later, on top, under its own
obligations.

\section{Verification of the Idempotence Theorem}
\label{ch05:sec:evidence}

\subsection{The exhaustive sweep and its count}

The executable suite verifies Theorem~\ref{ch05:thm:factorial} by complete
enumeration over small carriers: for every endomap $f$ on a representative
object of each size $N \in \{0, 1, 2, 3, 4\}$, it computes
$g = f^{N!}$ via the implemented exponentiation and asserts
$g \circ g = g$ by extensional arrow equality. The population, as computed
in Section~\ref{ch05:sec:construction}, is
\[
  \sum_{k=0}^{4} k^{k}
  = 0^{0} + 1^{1} + 2^{2} + 3^{3} + 4^{4}
  = 1 + 1 + 4 + 27 + 256
  = 289
\]
endomaps, matching the recorded suite count. The factorial exponents
involved are $0! = 1! = 1$, $2! = 2$, $3! = 6$, and $4! = 24$; the sweep
therefore exercises the exponentiation path with genuinely composite
binary expansions (e.g.\ $24 = 11000_2$) as well as the trivial ones.

Size four is a meaningful boundary for this particular theorem, not merely
an affordable one: carriers of size at most two cannot exhibit a
simultaneous nontrivial tail and nontrivial cycle, and size three barely
can; at size four the sweep includes maps with tail length up to three,
cycle lengths $1$ through $4$, mixed components, and every combination of
index and period realizable on four points. In small-scope terms
\cite{jackson}, the structural features the proof relies on --- tails,
cycles, divisibility of the exponent by the cycle length --- are all
present in force within the swept scope.

\begin{boundary}
\label{ch05:bound:testing}
The $289$-endomap sweep establishes the idempotence identity for every
endomap on carriers of size at most four, under the implemented
exponentiation, and nothing more. Theorem~\ref{ch05:thm:factorial} itself
holds for all finite sizes by the proofs above, which are standard
mathematics \cite{pin, howie}; but the connection between the theorem and
the running code is tested only within the stated bound, and the theorem
has no Lean formalization: ledger obligation PO-TKS8-CORE-04 stands at
\textsc{standard result / formalization open}. In particular, no artifact
should describe the factorial idempotence theorem as machine-verified.
\end{boundary}

\subsection{What closing the formalization obligation would involve}
\label{ch05:sec:po04}

Since PO-TKS8-CORE-04 is the most prominent open formalization in the
kernel's ledger, it is worth being concrete about its shape, so the
obligation reads as a bounded engineering task rather than an indefinite
aspiration. The mathematical content of
Theorem~\ref{ch05:thm:factorial} is elementary: the proof uses the
pigeonhole principle (Lemma~\ref{ch05:lem:repeat}), induction along
orbits, and the divisibility facts $\ell \le N \Rightarrow \ell \mid N!$
and $N! \ge N - 1$. None of this strains a proof assistant; comparable
statements about iterates of functions on finite types are routine in
mature libraries. The cost driver is the seed's own discipline: the
formal development is deliberately dependency-minimal and axiom-audited
(\texttt{propext} only, no external mathematics library), so the
formalization must either build the required elementary number theory
and finite-iteration lemmas in-house under the same audit, or extend the
audit policy to admit a vetted library dependency --- a policy decision
with consequences beyond this one theorem, since it would also move the
open norm-theorem and shortlex obligations. Either way, the theorem
statement to be formalized is exactly the extensional identity
$f^{N!} \circ f^{N!} = f^{N!}$ over the Lean model's arrows, stated for
every carrier size --- the model-level analogue of what the $289$-case
sweep checks for the Python artifact at sizes through four. Until that
is done, the theorem's status line remains what the ledger says:
standard result, formalization open.

\subsection{Why exhaustive rather than sampled}

For this theorem the exhaustive design is strictly superior to sampling at
equal cost, and it is worth recording why, since the reasoning generalizes.
The failure modes worth fearing are not uniform over the space of endomaps:
an exponentiation bug (say, a wrong squaring order) might affect only
exponents with a particular bit pattern; an off-by-one in the safe exponent
would spare permutations (whose idempotent power tolerates slack in the
index condition trivially when $i=0$) and bite only maps whose index is
maximal. Such structured defects occupy thin slices of the space, exactly
where uniform random sampling is weakest. Complete enumeration through size
four costs a few hundred cases --- trivially affordable --- and guarantees
that every realizable small configuration of tail, cycle, and exponent bit
pattern is hit. The suite's marginal design question is never ``how many
random cases''; it is ``what is the largest scope we can afford to
saturate,'' which is the small-scope methodology applied deliberately
\cite{jackson}.

\section{The Power API Contract}
\label{ch05:sec:apicontract}

Before descending into the exponentiation algorithm, we state the
behavioral contract of the power operation as an invariant, parallel to
the mismatch-rejection contract of Chapter~\ref{ch04:main}. The contract
is what the negative fixtures pin; the algorithm of the next section is
one implementation satisfying it.

\begin{invariant}[Power contract]
\label{ch05:inv:power}
The power operation is defined on certified endomaps and non-negative
integer exponents, and satisfies:
\begin{enumerate}
\item \textup{(Domain)} The operation exists only on the endomap wrapper
type; no arrow, path, or other value has it. The exponent must be a
non-negative integer; negative integers, non-integers, and the host
language's boolean values (which the host treats as integers, a coercion
the kernel explicitly refuses) are rejected with a validation error.
\item \textup{(Semantics)} The result is the certified endomap whose
arrow is $f^{n}$ in the sense of Definition~\ref{ch05:def:power}:
extensionally equal to the $n$-fold composite, with $f^{0} = \id_X$ on
the endomap's carrier.
\item \textup{(Closure)} The result is again a certified endomap on the
same carrier, so powers can be taken of powers, and
Proposition~\ref{ch05:prop:mult} governs the iteration.
\item \textup{(Name discipline)} The result's presentation name embeds
the bounded label of Construction~\ref{ch05:con:names}; the name has no
effect on the result's identity (Invariant~\ref{ch05:inv:names}).
\item \textup{(Resource bound)} Time is $O(|X| \log n)$ and space
$O(|X|)$, independent of $n$'s decimal size; in particular no operation
in the power path converts $n$ to decimal when $n$ exceeds the naming
threshold.
\end{enumerate}
\end{invariant}

\begin{checkclass}[Power fixtures]
\label{ch05:check:power}
The suite pins the contract with: rejection fixtures for exponent $-1$,
for a floating-point exponent, and for boolean exponents (both values);
the semantic fixtures $f^{0} = \id_X$, $f^{1} = f$, and agreement of
$f^{n}$ with the naive $n$-fold composite for all $n$ up to a small bound
on representative endomaps; the exhaustive idempotence sweep of
Section~\ref{ch05:sec:evidence}, which exercises exponents $1, 2, 6, 24$
across all $289$ small endomaps; and the large-exponent regressions at
$2^{30}$ and $10^{5000}$, which pin clauses (4) and (5). The
boolean-rejection fixture earns its place by a real hazard: in the host
language, \texttt{True} silently behaves as $1$ in arithmetic, so an
accidental boolean argument would otherwise compute $f^{1}$ without
complaint --- a plausible bug (a comparison result passed where a count
was meant) producing a plausible-looking wrong answer, which is the worst
combination.
\end{checkclass}

\begin{remark}[The idempotence predicate]
\label{ch05:rem:idempredicate}
Alongside \texttt{power}, the endomap type exposes a predicate testing
whether the map is idempotent, implemented as the direct comparison
$f \circ f = f$ --- one composition and one tuple equality, $O(|X|)$.
The predicate makes no structural analysis and returns no certificate
beyond its boolean; it exists so that the idempotence assertions of the
test suite are stated against the public API rather than against
internals, and so that downstream code can \emph{check} claimed
idempotents cheaply instead of trusting them. The pairing is
characteristic of the kernel's posture: the theorem
(Theorem~\ref{ch05:thm:factorial}) supplies an exponent that provably
works, and the predicate supplies an independent, trivially auditable
confirmation on each concrete instance.
\end{remark}

\section{Binary Exponentiation}
\label{ch05:sec:binexp}

\subsection{The algorithm}

Computing $f^{n}$ by $n - 1$ successive compositions is linear in $n$,
which is unusable at the exponents the kernel supports (the regression
suite goes to $10^{5000}$). The implementation uses the standard
square-and-multiply scheme, iterating over the binary expansion of the
exponent.

\begin{algorithm}[Exponentiation by squaring]
\label{ch05:alg:binexp}
Input: an endomap $f$ on $X$, an integer $n \ge 0$ (the implementation
validates that $n$ is a non-negative integer and explicitly rejects
booleans). Output: $f^{n}$.
\begin{enumerate}
\item Set $r \leftarrow \id_X$ (the accumulator), $b \leftarrow f$ (the
current square), $m \leftarrow n$ (the remaining exponent).
\item While $m > 0$:
  \begin{enumerate}
  \item if $m$ is odd, set $r \leftarrow b \circ r$;
  \item set $m \leftarrow \lfloor m / 2 \rfloor$;
  \item if $m > 0$, set $b \leftarrow b \circ b$.
  \end{enumerate}
\item Return the arrow with the accumulator's mapping tuple, re-wrapped as
an endomap on $X$ with a bounded presentation name
(Construction~\ref{ch05:con:names}).
\end{enumerate}
\end{algorithm}

\begin{proposition}[Correctness of Algorithm~\ref{ch05:alg:binexp}]
\label{ch05:prop:binexpcorrect}
Algorithm~\ref{ch05:alg:binexp} returns $f^{n}$.
\end{proposition}

\begin{proof}
We claim the loop maintains the invariant $r \circ$-combined with the
remaining work equals the answer; precisely:
\[
  b^{\,m} \circ r = f^{n}
\]
at every loop-condition evaluation, where $b^{m}$ denotes the $m$-th power
of the current square. Initialization: $b = f$, $m = n$, $r = \id_X$, and
$f^{n} \circ \id_X = f^{n}$ by right neutrality. Preservation: write
$m = 2q + \epsilon$ with $\epsilon \in \{0, 1\}$. If $\epsilon = 1$ the
body performs $r' = b \circ r$, and
\[
  (b \circ b)^{q} \circ r'
  = b^{2q} \circ (b \circ r)
  = b^{2q + 1} \circ r
  = b^{m} \circ r = f^{n},
\]
using Proposition~\ref{ch05:prop:mult} for $(b \circ b)^{q} = b^{2q}$,
associativity, and Proposition~\ref{ch05:prop:add}. If $\epsilon = 0$ the
same computation without the extra factor gives
$(b \circ b)^{q} \circ r = b^{m} \circ r = f^{n}$. (When $q = 0$ the
final squaring is skipped; the invariant holds vacuously with $b^{0} =
\id_X$.) Termination: $m$ strictly decreases toward $0$. At exit $m = 0$,
so $f^{n} = b^{0} \circ r = \id_X \circ r = r$ by left neutrality.
\end{proof}

\begin{example}[Trace at $n = 24$]
\label{ch05:ex:trace}
The safe exponent for a four-point carrier is $4! = 24 = 11000_2$, and the
algorithm's run at $n = 24$ is short enough to trace in full. Write the
state as $(m, r, b)$ at each loop test, with $r$ and $b$ given as powers
of $f$:
\[
\begin{array}{llll}
\text{iteration} & m & r & b \\
\text{start} & 24 & f^{0} & f^{1} \\
1 & 24 \text{ even; halve, square} & f^{0} & f^{2} \\
2 & 12 \text{ even; halve, square} & f^{0} & f^{4} \\
3 & 6 \text{ even; halve, square} & f^{0} & f^{8} \\
4 & 3 \text{ odd; multiply, halve, square} & f^{8} & f^{16} \\
5 & 1 \text{ odd; multiply, halve; final squaring skipped} & f^{24} & - \\
\end{array}
\]
Five iterations, six compositions (three squarings before the first
multiply, one more squaring, two multiplies), against twenty-three
compositions for the naive method. The invariant
$b^{m} \circ r = f^{24}$ of
Proposition~\ref{ch05:prop:binexpcorrect} can be checked at each row:
$f^{1 \cdot 24} \circ f^{0}$, $f^{2 \cdot 12} \circ f^{0}$,
$f^{4 \cdot 6} \circ f^{0}$, $f^{8 \cdot 3} \circ f^{0}$,
$f^{16 \cdot 1} \circ f^{8}$, and finally $f^{0} \circ f^{24}$ --- all
equal to $f^{24}$ by Proposition~\ref{ch05:prop:add}. Applied to the
endomap of Example~\ref{ch05:ex:mixed}, the returned arrow's tuple is
$(2, 1, 2, 1)$, the retraction onto the $2$-cycle, and its label is the
plain decimal \texttt{24}, far below the $256$-bit naming threshold of
Construction~\ref{ch05:con:names}.
\end{example}

\begin{remark}[Cost]
\label{ch05:rem:cost}
The loop executes $\lceil \log_2(n + 1) \rceil$ iterations, each performing
at most two compositions, and each composition of endo-arrows on $X$ is
$\Theta(|X|)$ tuple indexing. Total: $O(|X| \log n)$ time and $O(|X|)$
space per power. For $n = 10^{5000}$ this is roughly $1.7 \times 10^{4}$
iterations --- the exponent's bit length is about $16{,}610$ --- times a
four-entry tuple map in the regression's configuration: microseconds, not
megayears. The point of the $10^{5000}$ regression is not performance
bragging; it is that \emph{nothing} in the power path scales with the
exponent's decimal magnitude except the loop count, and the next
construction is what makes that true for the name as well.
\end{remark}

\section{Bounded Presentation Names}
\label{ch05:sec:names}

\subsection{The problem}

Every arrow carries a presentation name --- a human-readable string used in
error messages, logs, and debugging output. The natural naming scheme for a
power, something like \texttt{f\^{}(24)}, embeds the exponent in decimal.
For $n = 10^{5000}$ that name would be five thousand and one characters of
digits; worse, converting truly enormous integers to decimal is expensive
(quadratic in the classical algorithm), and the host language, as a
denial-of-service mitigation, enforces a hard limit on integer-to-decimal
conversion size, so a naive \texttt{str(n)} on a big-enough exponent does
not merely produce an unwieldy name --- it raises an exception and takes
the power operation down with it. A correctness-irrelevant string must
never be able to do that.

\subsection{The construction}

\begin{construction}[Bounded integer labels]
\label{ch05:con:names}
The name of a computed power $f^{n}$ embeds a \emph{bounded label} for $n$,
defined as follows:
\begin{enumerate}
\item if the bit length of $n$ is at most $256$, the label is the ordinary
decimal rendering of $n$;
\item otherwise, the label is the string
\[
  \texttt{bits=}b\texttt{,sha256=}d,
\]
where $b$ is the bit length of $n$ (a small integer) and $d$ is the first
sixteen hexadecimal characters of the SHA-256 digest of $n$'s big-endian
byte representation.
\end{enumerate}
In addition, the \emph{intermediate} arrows created inside
Algorithm~\ref{ch05:alg:binexp} --- the accumulator and the running square
--- carry fixed constant names, not compounded ones, so that name length
cannot grow with the iteration count; only the final returned arrow
receives the labeled name.
\end{construction}

The threshold and the digest length are engineering choices with simple
rationales. A $256$-bit threshold keeps every humanly-plausible exponent
(anything up to about $77$ decimal digits) fully readable, while
guaranteeing that decimal conversion is only ever invoked on integers far
below the host's conversion limit. Above threshold, the label's cost is
$O(\log n)$ --- dominated by the byte serialization of the exponent ---
and its size is constant. The sixteen-hex-character (64-bit) digest prefix
is not a cryptographic commitment; it is a collision-\emph{unlikely}
mnemonic that lets a human distinguish two large exponents in a log with
high probability. Nothing anywhere depends on the label being
collision-free, which is the subject of the next section.

The two regression pins bracket the construction's branches:
$n = 2^{30}$ has bit length $31 \le 256$ and must render in full decimal
(\texttt{1073741824}); $n = 10^{5000}$ has bit length about $16{,}610$ and
must render in the \texttt{bits=...,sha256=...} form, with the whole power
computation completing without touching large-integer decimal conversion.
Both cases assert, additionally, the mathematical correctness of the
resulting arrow --- the label machinery is exercised \emph{and} the
underlying mapping tuple is checked.

\subsection{Why names can never affect correctness}

\begin{invariant}[Extensional inertness of names]
\label{ch05:inv:names}
Presentation names participate in no semantic operation of the kernel.
Arrow equality and hashing exclude the name field entirely
(Definition~\ref{ch04:def:extensional}); the denotation key of an arrow is
the triple of source, target, and mapping tuple; composition, power,
idempotence testing, path composition, lift construction, and normal-form
computation read only denotational data. Consequently: two arrows with
colliding labels but different mapping data are unequal; two arrows with
different labels but equal denotational data are equal; and any name
whatsoever --- truncated, hashed, or constant --- yields the same results
in every kernel computation.
\end{invariant}

The invariant is enforced structurally rather than by policy: the equality
and hash implementations are generated from the dataclass declaration with
the name field marked non-comparing, so there is no code path in which a
name could be consulted, rather than a convention that no code path should.
The negative-style test for this is the pleasant one: construct the same
power twice under different naming (e.g.\ once through
Algorithm~\ref{ch05:alg:binexp} with its labeled name, once by repeated
composition with default composite names) and assert the results
\emph{equal}; construct two different arrows and force their names equal,
and assert them \emph{unequal}. Both directions are in the suite.

This is also why the truncated digest is harmless. A collision between two
labels --- possible in principle, vanishingly unlikely at 64 bits within
any log a human reads --- could at worst cause a moment of confusion in a
debugging session, at which point the denotation key, which is exact,
resolves the question. No proof, no test oracle, and no canonical form
consumes a label.

\begin{boundary}
\label{ch05:bound:names}
The naming construction makes a bounded-cost, bounded-size claim about
label \emph{construction}, verified by the two regression pins; it makes no
uniqueness or injectivity claim about labels, and no kernel artifact may
treat presentation names as identifiers. The unique identifier of an arrow
is its denotation key.
\end{boundary}

\section{Exercises}
\label{ch05:sec:exercises}

\begin{exercise}
\label{ch05:ex:analyze}
Analyze the endomap on five points with tuple $(1, 2, 3, 1, 0)$: draw the
functional graph, compute every tail length and cycle length, determine
$i$, $p$, and $k^{*}$ via Proposition~\ref{ch05:prop:whichpowers}, and
verify that $f^{5!} = f^{k^{*}}$ as Corollary~\ref{ch05:cor:whichidem}
predicts.
\end{exercise}

\begin{exercise}
\label{ch05:ex:extremal}
(i) Show that on $N$ points ($N \ge 2$) the index bound $i(f) \le N - 1$
is attained exactly by the maps whose functional graph is a single path of
$N - 1$ tail edges ending in a fixed point, and that those maps have
$p = 1$: maximal index forces trivial period. (ii) Show that the
\emph{pointwise} bound $t(x) + \ell(x) \le N$ of
Lemma~\ref{ch05:lem:enters} does \emph{not} globalize: the inequality
$i(f) + p(f) \le N$ fails for the permutation of five points with cycle
type $(2, 3)$, which has $i = 0$ and $p = \mathrm{lcm}(2, 3) = 6 > 5$.
Identify exactly where the globalization argument breaks: the period is an
lcm across components, while each point sees only its own cycle length.
\end{exercise}

\begin{exercise}
\label{ch05:ex:squares}
Enumerate the squares among the $27$ endomaps on three points: compute
$\{ g^{2} : g \in \Hom(X, X) \}$ for $|X| = 3$ and verify that the
$3$-cycles are squares (each $3$-cycle is the square of the other) while
the transpositions are not. Compare with the two-point enumeration in
Chapter~\ref{ch06:main}, where no nonidentity permutation is a square.
\end{exercise}

\begin{exercise}
\label{ch05:ex:lcmsharp}
Compute $N!$, $L(N) = \mathrm{lcm}(1, \dots, N)$, and the maximal
$k^{*}(f)$ over all endomaps on $N$ points, for $N = 1, \dots, 6$, and
confirm that $L(N) < N!$ strictly from $N = 4$ onward (they agree through
$N = 3$). (For the maximal $k^{*}$,
combine a long tail with a cycle structure of large lcm and note the
competition for points.)
\end{exercise}

\begin{exercise}
\label{ch05:ex:namescollide}
Using Invariant~\ref{ch05:inv:names}, argue that if two distinct exponents
$n_1 \ne n_2$ had colliding bounded labels
(Construction~\ref{ch05:con:names}) \emph{and} $f^{n_1} = f^{n_2}$ as
arrows, no kernel observation could distinguish the two computations ---
and explain why this is correct behavior rather than a defect, by
identifying which datum (the denotation key) is the arrow's identity and
which (the label) is commentary.
\end{exercise}

\section{Summary of Claims and Their Status}
\label{ch05:sec:summary}

\begin{center}
\small
\begin{tabular}{@{}p{5.6cm}p{7.2cm}@{}}
\toprule
Claim & Status \\
\midrule
Power laws (Propositions~\ref{ch05:prop:add}, \ref{ch05:prop:mult}) &
Complete informal proofs from the category laws; standard. \\
Functional-graph structure; index/period; monogenic structure
(Propositions~\ref{ch05:prop:shape}, \ref{ch05:prop:powerseq},
\ref{ch05:prop:monogenic}) &
Standard mathematics \cite{howie, pin}; complete informal proofs given as
background; no kernel API computes index or period. \\
Factorial idempotence (Theorem~\ref{ch05:thm:factorial}) &
Standard result \cite{pin, howie}; two complete informal proofs; exhaustive
executable verification for all $289 = \sum_{k \le 4} k^{k}$ endomaps
through carrier size four; Lean formalization \textsc{open}
(PO-TKS8-CORE-04). \\
Sharper bounds ($k^{*}(f)$; $\mathrm{lcm}(1..N)$;
Propositions~\ref{ch05:prop:whichpowers}, \ref{ch05:prop:lcm}) &
Standard mathematics, presented for orientation; \emph{not implemented,
not tested, no artifact} (Finding~\ref{ch05:find:implemented}). \\
Binary exponentiation correctness
(Proposition~\ref{ch05:prop:binexpcorrect}) &
Complete informal invariant proof; exercised by the idempotence sweep and
the large-exponent regressions. \\
Bounded presentation names (Construction~\ref{ch05:con:names}) &
Executable validation: regression pins at $2^{30}$ and $10^{5000}$; names
extensionally inert by construction
(Invariant~\ref{ch05:inv:names}). \\
\bottomrule
\end{tabular}
\end{center}

The endomap is the licensed half of the kernel's iteration story: one
arrow, one object, a genuine monoid, and a theorem with teeth. The next
chapter is the unlicensed half --- sequences of arrows across different
objects, which compose but do not iterate --- and the analysis of exactly
what goes wrong when the two halves are confused.
